\documentclass[11pt]{article}
\usepackage[margin=1in]{geometry}
\usepackage{amsmath,amssymb,bm}
\usepackage{booktabs,array}
\usepackage{enumitem}
\usepackage{float}
\usepackage{xcolor}
\usepackage{microtype}
\usepackage[hidelinks]{hyperref}

\definecolor{accent}{RGB}{35,65,105}
\definecolor{shade}{RGB}{245,247,250}

\title{\textbf{Steady Runaway Probability Function}\\[2mm]
Adjoint PDE, Finite-Volume Reference Solver, and FV-Assisted Parametric PINN}
\author{}
\date{September 2026}

\newcommand{\p}{p}
\newcommand{\x}{\xi}
\newcommand{\g}{\gamma}
\newcommand{\N}{N}
\newcommand{\f}{f}
\newcommand{\Pvec}{\bm P}
\newcommand{\Nvec}{\bm N}
\newcommand{\rRE}{\bm r_{\rm RE}}
\newcommand{\rF}{\bm r_{\rm F}}
\newcommand{\lam}{\bm\lambda}
\newcommand{\thetaNN}{\bm\theta}
\newcommand{\Ecfg}{\mathcal E}

\begin{document}
\maketitle

\begin{center}
\fcolorbox{accent}{shade}{\parbox{0.91\textwidth}{
\textbf{Purpose.} This document is the task-level scientific and numerical reference for the
steady runaway probability function (RPF).  It is intentionally narrower than the full
0D--2P kinetic-model reference.  The central objects are: (i) the steady first-passage
adjoint PDE, (ii) the conservative finite-volume (FV) discrete adjoint used as the numerical
reference, and (iii) the FV-assisted parametric PINN trained to represent the resulting RPF
family.  The FV reference, not the PINN, defines the qualified discrete ground truth for this
task.
}}
\end{center}

\medskip

The RPF asks a first-passage question for a tagged electron evolving under the local
small-angle kinetic generator: starting at $(p,\xi)$, what is the probability that the
electron exits the represented suprathermal domain through the successful runaway part of
the outer momentum boundary before it is lost through the lower-momentum failure boundary?
The answer is a bounded response function
\begin{equation}
 0\le P(p,\xi;\lam)\le 1,
\end{equation}
where $\lam$ denotes the prescribed plasma state.

The numerical hierarchy used throughout this document is
\begin{equation}
 \boxed{
 \begin{gathered}
 \text{local forward generator}
 \;\longrightarrow\;
 \text{exact FV transpose and first-passage solve}\\
 \longrightarrow\;\text{qualified FV labels}
 \;\longrightarrow\;
 \text{parametric PINN}.
 \end{gathered}
 }
\end{equation}
The PINN therefore accelerates evaluation of a model already established by the adjoint FV
solver; it does not introduce an independent definition of runaway probability.

\tableofcontents
\newpage

%===============================================================================
\section{Scope and model hierarchy}
\label{sec:scope}

This task concerns the \emph{steady} first-passage adjoint.  The following pieces are in
scope:
\begin{itemize}[leftmargin=2em]
 \item a spatially homogeneous, gyrotropic $(p,\xi)$ local electron generator;
 \item electric acceleration, asymptotically matched small-angle test-particle collisions,
       synchrotron reaction, partial screening, and the associated open radial boundaries;
 \item a conservative FV forward generator $L_0$ and its exact algebraic transpose;
 \item the steady discrete RPF solve $(-L_0^T)\Pvec=\rRE$;
 \item a parametric neural approximation $P_{\thetaNN}(p,\xi;\lam)$ trained with both the
       continuum adjoint residual and qualified FV labels;
 \item coefficient parity, FV refinement, held-out parameter validation, and optimizer/
       collocation refinement as separate qualification axes.
\end{itemize}

The following pieces are deliberately outside the present task-level document:
\begin{itemize}[leftmargin=2em]
 \item the full forward time-dependent distribution-function program;
 \item the finite-horizon time-dependent adjoint;
 \item self-consistent 0D plasma evolution;
 \item a fully discrete reverse-mode adjoint of a time-marching trajectory;
 \item differentiation through FV assembly or the sparse direct solver.
\end{itemize}

The Chiu--Harvey large-angle avalanche operator is also \emph{not} part of the tagged-
electron first-passage generator.  It is a branching source that creates additional
electrons.  Once an RPF has been computed, it may be contracted with such a source as a
response weight, but the branching process is not inserted into $L_0$ when defining the
single-particle RPF.

%===============================================================================
\section{Continuous steady first-passage problem}
\label{sec:continuous_rpf}

\subsection{Phase space and local forward generator}

Use the dimensionless relativistic momentum
\begin{equation}
 p=\frac{P}{m_ec},\qquad \gamma=\sqrt{1+p^2},\qquad -1\le\xi\le1.
\end{equation}
The represented RPF domain is
\begin{equation}
p_{\min}\le p\le p_{\max},\qquad -1\le\xi\le1,
\end{equation}
with $p_{\min}\ge0$.  The FV formulation admits the coordinate origin: its
first cell center remains strictly positive, coordinate-singular collision
coefficients are not evaluated on the $p=0$ face, and the radial face measure
vanishes there.

The local forward kinetic equation, with branching sources removed, is written in
conservative form as
\begin{equation}
 \frac{\partial f}{\partial\tau}
 =-\frac{1}{p^2}\frac{\partial}{\partial p}\left(p^2J_p\right)
  -\frac{\partial J_\xi}{\partial\xi}
 \equiv \mathcal L_0 f,
 \label{eq:forward_local}
\end{equation}
with
\begin{align}
 J_p&=A_p f-C_A\frac{\partial f}{\partial p},\\
 J_\xi&=A_\xi f-D_\xi\frac{\partial f}{\partial\xi},
\end{align}
and
\begin{align}
 A_p
 &=-\bar E\,\xi-\alpha\gamma p(1-\xi^2)-C_F,
 \label{eq:Ap}\\
 A_\xi
 &=(1-\xi^2)\left(-\frac{\bar E}{p}+\alpha\frac{\xi}{\gamma}\right),
 \label{eq:Axi}\\
 D_\xi
 &=\frac{\nu_D}{2}(1-\xi^2).
 \label{eq:Dxi}
\end{align}
Here $C_F$ is collisional radial drag, $C_A$ is radial energy diffusion, $\nu_D$ is the
pitch-angle deflection frequency, and $\alpha$ is the normalized synchrotron strength.
The coefficient construction uses only the asymptotically matched test-particle collision
operator, with energy-dependent Coulomb logarithms and partial-screening corrections.
Partial screening is always retained.
No separate adjoint collision model is introduced.

The parametric PINN configuration supplies the electric field as
\begin{equation}
 \Ecfg\equiv \frac{|E_\parallel|}{E_c},
\end{equation}
and the common physics layer maps the configured plasma state to the normalized operator
coefficients used in Eqs.~\eqref{eq:Ap}--\eqref{eq:Dxi}.  This notation separates the
user-facing parameter coordinate $\Ecfg$ from any lower-level reference normalization used
inside the kinetic implementation.

\subsection{Weighted adjoint PDE}

The physical momentum-space pairing is
\begin{equation}
 \langle P,f\rangle_\mu
 =2\pi\int_{p_{\min}}^{p_{\max}}p^2\,dp
      \int_{-1}^{1}P(p,\xi)f(p,\xi)\,d\xi.
 \label{eq:pairing}
\end{equation}
Away from the radial boundaries, integration by parts gives the backward operator
\begin{equation}
 \boxed{
 \mathcal L_0^\dagger P
 =A_p\frac{\partial P}{\partial p}
 +\frac{1}{p^2}\frac{\partial}{\partial p}
  \left(p^2C_A\frac{\partial P}{\partial p}\right)
 +A_\xi\frac{\partial P}{\partial\xi}
 +\frac{\partial}{\partial\xi}
  \left(D_\xi\frac{\partial P}{\partial\xi}\right).
 }
 \label{eq:continuous_adjoint}
\end{equation}
The steady interior RPF satisfies
\begin{equation}
 \mathcal L_0^\dagger P=0.
 \label{eq:continuous_steady}
\end{equation}
Equation~\eqref{eq:continuous_adjoint} is the continuum physics equation used by the PINN.
The production FV RPF is \emph{not} obtained by independently discretizing this equation;
it is obtained from the exact algebraic transpose of the conservative forward FV generator.
That distinction is central to the reference-solver architecture.

\subsection{First-passage boundary semantics}
\label{sec:boundary_semantics}

At the lower radial boundary $p=p_{\min}$, outward loss is defined as failure.  For the RPF,
\begin{equation}
 P(p_{\min},\xi)=0.
 \label{eq:lower_rpf_bc}
\end{equation}

The outer boundary is zero-inflow/open-outflow.  Its successful runaway subset is defined
from the same radial drift used by the forward operator,
\begin{equation}
 \Gamma_{\rm RE}(\lam)
 =\left\{(p_{\max},\xi):A_p(p_{\max},\xi;\lam)>0\right\}.
 \label{eq:success_set}
\end{equation}
On that successful subset, the first-passage probability is unity,
\begin{equation}
 P(p_{\max},\xi)=1,\qquad (p_{\max},\xi)\in\Gamma_{\rm RE}.
 \label{eq:upper_rpf_bc}
\end{equation}
If $\Gamma_{\rm RE}$ is empty, the successful-escape source vanishes naturally; no
hard-coded electric-field threshold is introduced.

At $\xi=\pm1$, both pitch drift and pitch diffusion vanish through the factor $1-\xi^2$,
so the pitch boundaries are natural zero-flux boundaries.

For the current FV-assisted PINN experiments, radial energy diffusion is disabled,
\begin{equation}
 C_A=0.
 \label{eq:no_energy_diffusion}
\end{equation}
The FV reference solver retains the more general capability, but the current neural
residual is intentionally restricted to Eq.~\eqref{eq:no_energy_diffusion}.

%===============================================================================
\section{Prescribed plasma parameterization}
\label{sec:parameters}

The current parametric surrogate uses deuterium plus one configurable impurity species.
The six-dimensional plasma parameter vector is
\begin{equation}
 \boxed{
 \lam=\left(\Ecfg,T_e,n_D,n_I,\overline Z_D,\overline Z_I\right).
 }
 \label{eq:lambda}
\end{equation}
Together with phase space, the network has eight inputs.

\subsection{Continuous mean-charge representation}

A noninteger mean charge is interpreted as a population mixture of neighboring integer
charge states \emph{before} collision coefficients are evaluated.  For species $s$ with
nuclear charge $Z_s$, density $n_s$, and mean charge $\overline Z_s\in[0,Z_s]$, define
\begin{equation}
 w_{s,q}(\overline Z_s)=\max\left(0,1-|\overline Z_s-q|\right),
 \qquad q=0,\ldots,Z_s,
 \label{eq:charge_weights}
\end{equation}
and
\begin{equation}
 n_{s,q}=n_s w_{s,q}.
\end{equation}
Then
\begin{equation}
 \sum_qn_{s,q}=n_s,
 \qquad
 \sum_q q n_{s,q}=n_s\overline Z_s.
\end{equation}
Quasineutrality and effective charge are therefore
\begin{align}
 n_e&=\sum_{s,q}q n_{s,q}=\sum_s n_s\overline Z_s,
 \label{eq:ne}\\
 Z_{\rm eff}&=\frac{1}{n_e}\sum_{s,q}q^2 n_{s,q}.
 \label{eq:zeff}
\end{align}
Partial-screening sums use the same resolved integer-state populations.  The FV reference
and the PINN coefficient layer must perform this same population-first mapping.  Final
collision coefficients are not interpolated between charge states.

\subsection{Normalized neural coordinates}

The network input is
\begin{equation}
 \bm z=(\widehat p,\widehat\xi,\widehat\lam)\in[0,1]^8.
\end{equation}
Momentum is logarithmically normalized,
\begin{equation}
 \widehat p
 =\frac{\ln p-\ln p_{\min}}
        {\ln p_{\max}-\ln p_{\min}},
 \qquad
 \widehat\xi=\frac{\xi+1}{2}.
 \label{eq:phase_norm}
\end{equation}
The electric field, temperature, and densities may be mapped linearly or logarithmically
according to the training configuration; mean charge states are mapped linearly.  The
physical surrogate domain is therefore a configuration choice rather than a hard-coded
network property.

%===============================================================================
\section{Finite-volume reference solver}
\label{sec:fv_reference}

The FV solver is the numerical reference for the steady RPF.  It first constructs the
qualified local forward generator and then forms the exact discrete transpose.  The
reference is therefore a \emph{discretize-then-adjoint} construction.

\subsection{Cell-content state and conservative grid}

The baseline grid is a tensor product with
\begin{equation}
 \Delta p=\frac{p_{\max}-p_{\min}}{N_p},\qquad
 \Delta\xi=\frac{2}{N_\xi},
\end{equation}
cell centers
\begin{equation}
 p_i=p_{\min}+\left(i-\frac12\right)\Delta p,
 \qquad
 \xi_j=-1+\left(j-\frac12\right)\Delta\xi,
\end{equation}
and corresponding face locations.  The primary unknown is the physical cell content
\begin{equation}
 N_{ij}
 =2\pi\int_{p_{i-1/2}}^{p_{i+1/2}}p^2\,dp
      \int_{\xi_{j-1/2}}^{\xi_{j+1/2}}f(p,\xi)\,d\xi.
 \label{eq:cell_content}
\end{equation}
The exact cell phase volume is
\begin{equation}
 V_{ij}=\frac{2\pi}{3}\left(p_{i+1/2}^3-p_{i-1/2}^3\right)\Delta\xi,
 \qquad f_{ij}=\frac{N_{ij}}{V_{ij}}.
\end{equation}
Because the phase-space measure is already included in $\bm N$, the discrete dual pairing
is simply
\begin{equation}
 \langle\Pvec,\Nvec\rangle=\Pvec^T\Nvec.
 \label{eq:discrete_pairing}
\end{equation}
There is no additional FV mass matrix in the adjoint problem.

\subsection{Conservative local fluxes}

Integrating Eq.~\eqref{eq:forward_local} over a cell gives
\begin{equation}
 \frac{dN_{ij}}{d\tau}
 =-\left(F^p_{i+1/2,j}-F^p_{i-1/2,j}\right)
  -\left(F^\xi_{i,j+1/2}-F^\xi_{i,j-1/2}\right).
 \label{eq:fv_balance}
\end{equation}
Each interior face flux is formed once and enters the two neighboring cells with opposite
signs, so interior conservation is algebraic.

Both drift--diffusion fluxes use Chang--Cooper exponential fitting.  For a face in a
coordinate $q$ with drift $A_q$, diffusion coefficient $D_q$, and spacing $\Delta q$, define
\begin{equation}
 w_q=\frac{A_q\Delta q}{D_q},\qquad
 \delta(w_q)=\frac{1}{w_q}-\frac{1}{e^{w_q}-1},\qquad
 \delta(0)=\frac12,
 \label{eq:cc_weight}
\end{equation}
and the face value
\begin{equation}
 f^{\rm CC}_{q+1/2}
 =(1-\delta(w_q))f_q+\delta(w_q)f_{q+1}.
 \label{eq:cc_face_value}
\end{equation}
This orientation matches the flux convention $J_q=A_qf-D_q\partial_qf$ and preserves
a locally constant-coefficient zero-flux exponential.  When $D_q\to0$, the continuous
Chang--Cooper limit supplies the corresponding donor-side value.

For example, at an interior radial face,
\begin{equation}
 F^p_{i+1/2,j}
 =2\pi p_{i+1/2}^2\Delta\xi
 \left[
  A^p_{i+1/2,j}f^{p,\mathrm{CC}}_{i+1/2,j}
  -C_{A,i+1/2}\frac{f_{i+1,j}-f_{ij}}{\Delta p}
 \right].
 \label{eq:radial_flux}
\end{equation}
The pitch-face flux is
\begin{equation}
 F^\xi_{i,j+1/2}
 =V^p_i
  \left[
   A^\xi_{i,j+1/2}f^{\rm CC}_{i,j+1/2}
   -D_{\xi,i,j+1/2}\frac{f_{i,j+1}-f_{ij}}{\Delta\xi}
  \right],
 \qquad
 V^p_i=\frac{2\pi}{3}\left(p_{i+1/2}^3-p_{i-1/2}^3\right).
 \label{eq:pitch_flux}
\end{equation}

At $p=p_{\max}$ the open outer flux is
\begin{equation}
 F^p_{N_p+1/2,j}
 =2\pi p_{\max}^2\Delta\xi
  \left[A_p(p_{\max},\xi_j)\right]_+f_{N_p,j},
 \label{eq:outer_flux}
\end{equation}
with no incoming exterior distribution.  At the shifted inner boundary, outward loss is
absorbing failure.  These same boundary terms define the first-passage source vectors used
below.

\subsection{Discrete generator and exact transpose}

With branching sources disabled, the cell-content system is
\begin{equation}
 \boxed{
 \frac{d\Nvec}{d\tau}=L_0\Nvec.
 }
 \label{eq:L0}
\end{equation}
The steady discrete adjoint matrix is
\begin{equation}
 A\equiv -L_0^T.
 \label{eq:A_matrix}
\end{equation}
The implementation does not derive a second FV stencil for the adjoint.  It assembles the
qualified forward values of $-L_0$ and transposes those values exactly on the same
five-point CSR graph.

The successful outer-escape rate vector is defined by
\begin{equation}
 \Phi_{\rm RE}^{\rm FV}(\Nvec)=\rRE^T\Nvec,
\end{equation}
where, for the zero-inflow outer flux in the present $C_A=0$ model,
\begin{equation}
 (r_{\rm RE})_{ij}
 =\begin{cases}
 0, & i<N_p,\\[2mm]
 \displaystyle
 \frac{2\pi p_{\max}^2\Delta\xi}{V_{N_pj}}
 \left[A_p(p_{\max},\xi_j)\right]_+,
 & i=N_p.
 \end{cases}
 \label{eq:rre}
\end{equation}
The lower absorbing boundary similarly defines the failure vector $\rF$.

\subsection{Steady first-passage linear system}

If the local forward state decays at long time, then the integrated successful escape from
an arbitrary initial cell-content state is represented by
\begin{equation}
 \int_0^\infty \rRE^T e^{L_0\tau}\Nvec^0\,d\tau
 =\Pvec^T\Nvec^0,
\end{equation}
where
\begin{equation}
 \boxed{
 A\Pvec=\rRE,
 \qquad A=-L_0^T.
 }
 \label{eq:discrete_rpf}
\end{equation}
Thus $P_k$ is the eventual successful-escape probability of a unit cell-content state in
cell $k$.

The exact discrete duality identity is
\begin{equation}
 \boxed{
 \Pvec^T(-L_0)\Nvec=\rRE^T\Nvec
 }
 \label{eq:duality}
\end{equation}
for every test vector $\Nvec$ when Eq.~\eqref{eq:discrete_rpf} is solved exactly.

With no other sinks or sources, the total boundary-loss identity is
\begin{equation}
 -L_0^T\bm 1=\rRE+\rF,
 \label{eq:escape_identity}
\end{equation}
and the corresponding failure probability satisfies
\begin{equation}
 A\bm P_{\rm F}=\rF,
 \qquad
 \Pvec+\bm P_{\rm F}=\bm1.
 \label{eq:partition}
\end{equation}
The expected bounds are therefore $0\le P_k\le1$.

\subsection{GPU sparse-direct realization}

The production reference path uses a fixed five-point CSR topology.  Matrix values and
right-hand sides are assembled in FP64 on the GPU.  Warp-owned device arrays are passed
directly to cuDSS for symbolic analysis, numerical factorization, and sparse-direct
solution.  Symbolic analysis may be reused whenever the graph is unchanged and only
physical coefficients vary between parameter cases.

A normalized linear residual is evaluated from the actual GPU matrix and solution,
\begin{equation}
 \epsilon_{\rm solve}
 =\frac{\|A\Pvec-\rRE\|_2}{\|\rRE\|_2},
 \label{eq:linear_residual}
\end{equation}
with an absolute residual used when $\rRE=0$.

\subsection{Reference-solver qualification requirements}
\label{sec:fv_qualification}

An FV field is admitted as a PINN label only after the reference calculation passes the
relevant numerical checks.  At minimum:
\begin{enumerate}[leftmargin=2.3em]
 \item \textbf{Exact transpose:} the GPU values of $-L_0^T$ agree with an independently
       formed transpose of the assembled forward matrix to floating-point accuracy.
 \item \textbf{Successful-escape vector:} nonzero entries of $\rRE$ occur only in the
       outer shell and the pitch mask agrees with $A_p(p_{\max},\xi_j)>0$.
 \item \textbf{Sparse residual:} Eq.~\eqref{eq:linear_residual} is small relative to the
       spatial discretization error.
 \item \textbf{Discrete duality:} Eq.~\eqref{eq:duality} holds for deterministic and
       random nonnegative test states.
 \item \textbf{Probability partition:} the success/failure solutions satisfy
       Eq.~\eqref{eq:partition} and material violations of $0\le P\le1$ are absent.
 \item \textbf{Forward reconstruction:} direct propagation of selected initial states
       reproduces $\Pvec^T\Nvec^0$ as the integrated successful escape.
 \item \textbf{Grid and domain convergence:} $N_p$, $N_\xi$, $p_{\min}$, and $p_{\max}$
       are refined independently until the RPF quantities of interest are stable.
\end{enumerate}
A converged cuDSS residual alone does not establish a converged FV label.

\subsection{Wide momentum domains and conditioning of conservation diagnostics}
\label{sec:wide_domain}

A very broad interval $p_{\min}\ll1\ll p_{\max}$ is a difficult regime for a uniform
momentum mesh.  At low momentum the pitch-scattering coefficient becomes rapidly varying
and, in the suprathermal asymptotic regime, behaves schematically as
\begin{equation}
 \nu_D\sim p^{-3}.
\end{equation}
The discrete pitch-diffusion entries therefore scale roughly as
\begin{equation}
 \frac{\nu_D}{\Delta\xi^2}.
\end{equation}
Refining a uniform radial mesh moves the first cell center to smaller $p$, which can make
individual matrix coefficients much larger even while the desired first-passage response
remains bounded.

This matters for diagnostics such as Eq.~\eqref{eq:escape_identity}.  In exact arithmetic,
large interior contributions cancel and leave only boundary escape.  In floating-point
arithmetic, evaluating $A\bm1$ may therefore be a strongly cancellation-conditioned
operation.  A raw relative forward error can increase with refinement even when a
componentwise backward error remains near machine precision.  Such behavior must be
interpreted as a conditioning question, not automatically as either proof of correctness
or proof of failure.

Most importantly, good backward algebraic consistency does \emph{not} establish spatial
resolution.  Before broad-domain FV fields are promoted to PINN labels, the momentum grid
must resolve the low-$p$ region and the RPF itself must converge under grid/domain
refinement.  If a uniform mesh cannot do this economically, a nonuniform radial grid is a
numerical-model improvement to be qualified explicitly rather than hidden by relaxed
algebraic tolerances.

%===============================================================================
\section{FV-assisted parametric PINN}
\label{sec:pinn}

The PINN represents the steady RPF continuously in $(p,\xi)$ and over the six-dimensional
plasma parameter domain of Eq.~\eqref{eq:lambda}.  It is trained simultaneously against
the continuum adjoint equation and qualified FV solutions.

\subsection{Neural probability representation}

Let $\mathcal N_{\thetaNN}(\bm z)$ be the scalar output of a fully connected network.  The
current hard-output transform is
\begin{equation}
 \boxed{
 P_{\thetaNN}(\bm z)
 =\tanh\left[\widehat p^{\,2}\mathcal N_{\thetaNN}(\bm z)^2\right].
 }
 \label{eq:output_transform}
\end{equation}
It enforces
\begin{equation}
 P_{\thetaNN}(p_{\min},\xi;\lam)=0,
 \qquad
 0\le P_{\thetaNN}<1.
\end{equation}
The successful $p_{\max}$ condition is imposed through the boundary loss instead of being
hard-coded into the transform.

\subsection{Steady adjoint residual used for training}

For the present $C_A=0$ surrogate, Eq.~\eqref{eq:continuous_adjoint} becomes
\begin{align}
 \mathcal R_{\thetaNN}
 ={}&-U_p\frac{\partial P_{\thetaNN}}{\partial p}
 +(1-\xi^2)
  \left(\frac{\bar E}{p}-\frac{\alpha\xi}{\gamma}\right)
  \frac{\partial P_{\thetaNN}}{\partial\xi}
 \nonumber\\
 &-\frac{\nu_D}{2}\frac{\partial}{\partial\xi}
 \left[(1-\xi^2)\frac{\partial P_{\thetaNN}}{\partial\xi}\right],
 \label{eq:pinn_residual}
\end{align}
where
\begin{equation}
 U_p=A_p=-\bar E\xi-\alpha\gamma p(1-\xi^2)-C_F.
\end{equation}
Equivalently,
\begin{equation}
 \mathcal R_{\thetaNN}
 =a_pP_p+a_\xi P_\xi+a_{\xi\xi}P_{\xi\xi},
 \label{eq:expanded_residual}
\end{equation}
with
\begin{align}
 a_p&=-U_p,\\
 a_\xi&=(1-\xi^2)\left(\frac{\bar E}{p}-\frac{\alpha\xi}{\gamma}\right)+\nu_D\xi,\\
 a_{\xi\xi}&=-\frac12\nu_D(1-\xi^2).
 \label{eq:pinn_coeffs}
\end{align}

For a fixed collocation set, all physical quantities in
Eq.~\eqref{eq:pinn_coeffs} are independent of the network weights.  The implementation
therefore precomputes the normalized-coordinate versions of
$a_p$, $a_\xi$, and $a_{\xi\xi}$ once for each sampled PDE set.  Every optimizer step then
repeats only the neural differentiation required for $P_p$, $P_\xi$, and $P_{\xi\xi}$.
One first-derivative evaluation supplies both phase-space gradients and a directional JVP
supplies the required second pitch derivative.

\subsection{FV teacher data}

Plasma states $\lam_k$ are sampled over the configured six-dimensional domain.  For each
state, the standalone reference solver computes
\begin{equation}
 A(\lam_k)\Pvec_k^{\rm FV}=\rRE(\lam_k).
 \label{eq:teacher}
\end{equation}
The FV solve remains completely outside JAX automatic differentiation.  The dataset stores
parameter coordinates, FV phase-space coordinates, RPF fields, and qualification metadata.

Complete plasma states are split into training and validation subsets.  A configurable
subset of FV cells from each training state contributes to the supervised loss.  Held-out
parameter states are never used to update the network and are evaluated on their complete
FV grids after training.

\subsection{Composite training objective}

The loss is
\begin{equation}
 \boxed{
 \mathcal L(\thetaNN)
 =w_{\rm PDE}\mathcal L_{\rm PDE}
 +w_{\rm BC}\mathcal L_{\rm BC}
 +w_{\rm FV}\mathcal L_{\rm FV}.
 }
 \label{eq:total_loss}
\end{equation}
For fixed PDE collocation points,
\begin{equation}
 \mathcal L_{\rm PDE}
 =\frac{1}{N_{\rm PDE}}
  \sum_{i=1}^{N_{\rm PDE}}
  \mathcal R_{\thetaNN}(\bm z_i^{\rm PDE})^2.
\end{equation}
For high-$p$ boundary points, let $m_i=1$ on the successful set of
Eq.~\eqref{eq:success_set} and $m_i=0$ otherwise.  Then
\begin{equation}
 \mathcal L_{\rm BC}
 =\frac{\sum_i m_i[P_{\thetaNN}(\bm z_i^{\rm BC})-1]^2}
        {\max(\sum_i m_i,1)}.
\end{equation}
For supervised FV samples,
\begin{equation}
 \mathcal L_{\rm FV}
 =\frac{1}{N_{\rm FV}}
  \sum_{j=1}^{N_{\rm FV}}
  [P_{\thetaNN}(\bm z_j^{\rm FV})-P_j^{\rm FV}]^2.
\end{equation}
The sample counts and the three weights are configuration parameters.

\subsection{Whole-batch and chunked evaluation}

The mathematical objective is full batch.  To tune GPU performance and memory, each loss
sum may either be evaluated in one XLA call or accumulated internally in fixed-size chunks.
Chunking does not create stochastic minibatches and does not change the collocation set.
The configuration convention is
\begin{equation}
 \texttt{chunk}=0\quad\Longrightarrow\quad\text{evaluate the complete term at once}.
\end{equation}
A positive chunk size evaluates the same sum in deterministic chunks.  The choice is a
hardware/performance parameter and should be benchmarked on the target GPU.

%===============================================================================
\section{Optimization protocol}
\label{sec:optimization}

The training hierarchy is
\begin{equation}
 \boxed{
 \text{full-batch Adam warmup}
 \;\longrightarrow\;
 \text{blockwise full-batch SSBroyden refinement}.
 }
\end{equation}
This follows the high-accuracy PINN strategy in which a first-order optimizer reaches a
useful basin and a dense quasi-Newton method performs the final high-accuracy refinement
\cite{Urban2025PINN}.

\subsection{Adam stage}

Adam operates on one fixed PDE, boundary, and FV training set for the configured number of
steps.  The purpose of this stage is not necessarily to minimize each component of the
composite objective to high accuracy; it supplies the starting point for the subsequent
quasi-Newton phase.

\subsection{SSBroyden blocks and line search}

The default SSBroyden globalization is a strong-Wolfe backtracking line search.  The
parameters are flattened to one vector and the dense inverse-Hessian approximation is
maintained within each quasi-Newton solve.

The total SSBroyden phase is divided into configurable blocks.  Each block is a new solve,
so its inverse-Hessian approximation starts fresh.  Between blocks, PDE and outer-boundary
collocation points may optionally be resampled while FV labels remain fixed.  Thus the
training protocol can be written schematically as
\begin{equation}
 \text{SSB block}
 \;\to\;
 \left[\text{optional PDE/BC resample}\right]
 \;\to\;
 \left[\text{reset inverse Hessian}\right]
 \;\to\;
 \text{next SSB block}.
\end{equation}

This is a generalization-control mechanism rather than an arbitrary implementation detail.
A very accurate quasi-Newton method can overfit a finite collocation set when that set is
not sufficiently dense.  Resampling changes the discrete approximation to the continuum
physics objective, while resetting the inverse-Hessian approximation avoids carrying
secant information from the old sampled objective into the new one.  The resampling
cadence and block length are therefore problem-dependent training parameters and may be
disabled when a fixed collocation set is demonstrably sufficient.

Machine-precision optimizer tolerances apply only to the sampled finite-dimensional loss.
They are not evidence by themselves of continuum, FV-grid, or parametric-surrogate
convergence.

%===============================================================================
\section{Reference-to-surrogate workflow}
\label{sec:workflow}

The intended workflow is summarized in Table~\ref{tab:workflow}.

\begin{table}[htbp]
\centering
\small
\renewcommand{\arraystretch}{1.25}
\begin{tabular}{>{\raggedright\arraybackslash}p{0.12\textwidth}
                >{\raggedright\arraybackslash}p{0.25\textwidth}
                >{\raggedright\arraybackslash}p{0.52\textwidth}}
\toprule
\textbf{Stage} & \textbf{Object} & \textbf{Requirement}\\
\midrule
1 & Shared physics coefficients & FV and JAX/PINN coefficient layers agree, including
continuous mean-charge mapping, quasineutrality, screening, and synchrotron strength.\\
2 & FV operator & Assemble the conservative forward cell-content generator with the
selected boundaries and physical coefficients.\\
3 & Discrete adjoint & Form $A=-L_0^T$ by exact algebraic transpose and construct the
successful-escape vector from the same outer flux.\\
4 & FV reference solution & Solve $A\Pvec=\rRE$ and pass algebraic, probability, grid, and
domain qualification.\\
5 & Dataset & Store qualified RPF fields for training states and disjoint complete plasma
states for validation.\\
6 & PINN & Train the continuum adjoint residual, successful outer BC, and FV supervised
loss with Adam followed by SSBroyden.\\
7 & Surrogate qualification & Compare complete held-out FV fields, refine collocation and
FV label sets, and increase network capacity when representation error is limiting.\\
\bottomrule
\end{tabular}
\caption{Reference-to-surrogate hierarchy.  No stage downstream of the FV reference can
repair an unqualified reference discretization.}
\label{tab:workflow}
\end{table}

The practical separation of responsibilities is therefore:
\begin{center}
\fbox{\parbox{0.90\textwidth}{
\textbf{FV solver:} defines and qualifies the discrete reference RPF.\\[1mm]
\textbf{PINN:} approximates the qualified RPF family continuously in phase space and parameters.
}}
\end{center}

%===============================================================================
\section{Surrogate validation and promotion criteria}
\label{sec:surrogate_qualification}

Training loss alone is not a promotion criterion.  The minimum validation hierarchy is:
\begin{enumerate}[leftmargin=2.3em]
 \item \textbf{Coefficient parity.}  Compare the FV and JAX coefficient layers at
       representative points spanning the requested parameter domain.
 \item \textbf{FV reference convergence.}  Demonstrate RPF stability under $N_p$, $N_\xi$,
       $p_{\min}$, and $p_{\max}$ refinement before using a family of fields as labels.
 \item \textbf{Held-out plasma states.}  Exclude complete parameter states from training
       and compare the PINN with the complete FV field using at least relative $L^2$ and
       maximum absolute error.
 \item \textbf{Collocation/data refinement.}  Increase PDE points, boundary points, number
       of FV parameter cases, and FV samples per case until held-out quantities stabilize.
 \item \textbf{Optimization refinement.}  Vary Adam length, SSBroyden block length,
       resampling cadence, and optimizer tolerances as numerical-training convergence axes.
 \item \textbf{Network-capacity refinement.}  Increase width/depth when the FV error
       plateaus while PDE and boundary residuals are already small.
\end{enumerate}

A characteristic capacity-limited regime is one in which
\begin{equation}
 \mathcal L_{\rm PDE},\mathcal L_{\rm BC}\ll\mathcal L_{\rm FV},
\end{equation}
while held-out FV error no longer improves materially.  In that regime, further reduction
of the collocation residual is not the primary task; representation capacity, label
coverage, or parameter-domain complexity should be examined.

%===============================================================================
\section{Configuration quantities that belong to this task}
\label{sec:config}

The task-level configuration should expose, rather than hard-code, at least the following:

\begin{table}[H]
\centering
\small
\renewcommand{\arraystretch}{1.22}
\begin{tabular}{>{\raggedright\arraybackslash}p{0.24\textwidth}
                >{\raggedright\arraybackslash}p{0.66\textwidth}}
\toprule
\textbf{Category} & \textbf{User-configurable quantities}\\
\midrule
FV reference & $N_p$, $N_\xi$, $p_{\min}$, $p_{\max}$, magnetic field, energy-diffusion
switch, number of training and validation plasma states.\\
Parameter domain & lower/upper bounds and linear/log sampling transforms for $\Ecfg$, $T_e$,
$n_D$, $n_I$, together with bounds for $\overline Z_D$ and $\overline Z_I$.\\
PINN architecture & width, depth, initialization seed.\\
Training sets & $N_{\rm PDE}$, $N_{\rm BC}$, FV samples per training state, sampling seed,
interior-boundary exclusion.\\
Loss evaluation & PDE, BC, FV, and validation chunk sizes, including zero for whole-batch
evaluation.\\
Adam & number of steps, compiled block size, learning rate.\\
SSBroyden & number of blocks, iterations per block, resampling cadence, line-search choice,
relative and absolute tolerances.\\
Loss weights & $w_{\rm PDE}$, $w_{\rm BC}$, $w_{\rm FV}$ and any documented residual scale.\\
\bottomrule
\end{tabular}
\caption{Configuration contract for the focused RPF/FV/PINN task.}
\end{table}

The physical training domain and the training protocol are deliberately separate.  A
change in optimizer settings does not redefine the RPF, while expanding the physical
parameter or momentum domain can require new FV grid/domain qualification before new
labels are accepted.

%===============================================================================
\section{Summary of the task-level scientific contract}

The focused contract can be reduced to five statements:
\begin{enumerate}[leftmargin=2.3em]
 \item The RPF is a first-passage probability of the \emph{local} tagged-electron
       generator, with inner radial loss interpreted as failure and positive outer radial
       outflow interpreted as successful runaway.
 \item The authoritative numerical RPF is the solution of
       $(-L_0^T)\Pvec=\rRE$, where $L_0$ is the conservative FV forward cell-content
       generator and the transpose is algebraic and exact on the assembled graph.
 \item The FV reference must satisfy algebraic, probability, grid, and domain
       qualification before its fields are used as neural labels.
 \item The PINN solves the same steady adjoint physics continuously and is anchored by FV
       supervision.  It does not differentiate through the FV solver and does not replace
       the reference discretization.
 \item Adam--SSBroyden training, collocation resampling, Hessian resets, chunking, and
       network capacity are numerical training choices whose adequacy is established by
       held-out FV validation and refinement, not by training loss alone.
\end{enumerate}

%===============================================================================
\footnotesize
\begin{thebibliography}{99}

\bibitem{Hesslow2018Collision}
L.~Hesslow, O.~Embr\'eus, M.~Hoppe, T.~C. DuBois, G.~Papp,
M.~Rahm, and T.~F{\"u}l{\"o}p,
``Generalized collision operator for fast electrons interacting with partially ionized impurities,''
\emph{J. Plasma Phys.} \textbf{84}, 905840605 (2018).

\bibitem{Andersson2001}
F.~Andersson, P.~Helander, and L.-G.~Eriksson,
``Damping of relativistic electron beams by synchrotron radiation,''
\emph{Phys. Plasmas} \textbf{8}, 5221 (2001).

\bibitem{McDevitt2025Dynamics}
C.~J. McDevitt, J.~S. Arnaud, and X.-Z. Tang,
``A physics-constrained deep learning treatment of runaway electron dynamics,''
\emph{Phys. Plasmas} \textbf{32}, 042503 (2025).

\bibitem{Urban2025PINN}
J.~F. Urb{\'a}n, P.~Stefanou, and J.~A. Pons,
``Unveiling the optimization process of physics informed neural networks: How accurate and competitive can PINNs be?,''
\emph{J. Comput. Phys.} \textbf{523}, 113656 (2025).

\end{thebibliography}

\end{document}
